\chapter*{Introduction Générale}
\phantomsection
\addcontentsline{toc}{chapter}{Introduction Générale}

À l'ère de la quatrième révolution industrielle, le secteur de l'éducation connaît une mutation profonde portée par la convergence de l'intelligence artificielle et des technologies de collaboration en temps réel \cite{ai_education}. Cette transformation numérique ne se limite plus à la simple numérisation des supports pédagogiques, mais vise désormais à résoudre le défi majeur de l'enseignement moderne : passer d'un apprentissage passif et uniforme à une expérience active, immersive et personnalisée.

C'est dans ce contexte de rupture technologique que s'inscrit notre Projet de Fin d'Études intitulé \textbf{EduGenius}. Ce projet consiste en la conception et la réalisation d'une plateforme d'apprentissage de nouvelle génération qui fusionne les capacités de l'Intelligence Artificielle (IA) avec des outils de communication synchrone. L'ambition de cette solution est de redéfinir la relation entre l'apprenant, l'enseignant et le savoir en automatisant les tâches à faible valeur ajoutée pour se concentrer sur l'accompagnement pédagogique et la réussite académique.

La plateforme \textbf{EduGenius} est conçue comme un écosystème intégré répondant aux besoins de trois acteurs clés : l'administration, les enseignants et les étudiants. Elle se distingue par l'intégration native de moteurs d'IA pour la génération automatique de supports (résumés, quiz et flashcards), un module de visioconférence haute fidélité avec détection automatique de présence, et un système d'analyses avancées permettant de suivre proactivement les performances des apprenants. En s'appuyant sur des dynamiques de gamification, le projet vise non seulement à faciliter la transmission des connaissances, mais aussi à maximiser l'engagement et la rétention sur le long terme.

Le développement de cette solution a été mené selon une \textbf{démarche Agile (Scrum)}, garantissant une flexibilité maximale face aux défis techniques liés au traitement de données en temps réel et à l'intégration des modèles de langage. Cette méthodologie nous a permis de livrer des incréments fonctionnels réguliers tout en assurant une rigueur logicielle conforme aux standards de l'ingénierie moderne.

Le présent rapport documente l'ensemble du cycle de vie du projet, structuré autour des axes suivants :

\begin{itemize}
    \item \textbf{Chapitre 1 : Cadre général du projet} -- expose la problématique, les objectifs stratégiques, l'étude comparative des solutions du marché ainsi que le cadre méthodologique retenu.
    \item \textbf{Chapitre 2 : Analyse et Spécification des besoins} -- définit les exigences fonctionnelles et non-fonctionnelles à travers une modélisation rigoureuse des besoins utilisateurs.
    \item \textbf{Chapitre 3 : Conception de la solution} -- détaille l'architecture logicielle, la conception de la base de données et les choix de design système.
    \item \textbf{Chapitre 4 : Réalisation et Mise en œuvre} -- présente l'environnement technique, les technologies utilisées (Stack MERN) et les fonctionnalités clés implémentées.
    \item \textbf{Chapitre 5 : Tests et Évaluation} -- récapitule les phases de validation, les tests de performance et les résultats obtenus par rapport aux objectifs initiaux.
\end{itemize}